\section{Properties of Integrals}
%%%%%
\index{integrals!properties}%
\index{properties!of integrals}%

\begin{theorem}[additivity of the integral]
\mymark{*}%
\mymargin{additivity of the integral}%
\index{additivity of the integral}%
\index{integral!additivity}%
\label{th:additivita_integrale}%
Let $f\colon [a,b]\to \RR$ be a bounded function and let $c\in [a,b]$.
Then $f$ is Riemann-integrable on $[a,b]$ if and only if
$f$ is Riemann-integrable on $[a,c]$ and on $[c,b]$.
In such a case, it holds that
\begin{equation}\label{eq:36645}
  \int_a^b f = \int_a^c f + \int_c^b f.
\end{equation}

According to the convention
\[
    \int_b^a f = -\int_a^b f
\]
the formula \eqref{eq:36645} is valid not only if $a\le c\le b$ but also
if $a,b,c$ are in any order, provided that the function $f$ is integrable
on the interval containing all three points $a,b,c$.
\end{theorem}
%
\begin{proof}
\mymark{*}
Suppose that $f$ is integrable on $[a,c]$ and on $[c,b]$.
Then, based on the criteria for integrability, for every $\eps>0$ there exist a
partition $P$ of $[a,c]$ and a partition $Q$ of $[c,b]$ such that
\[
  S^*(f,P) - S_*(f,P) < \frac \eps 2,
  \qquad
  S^*(f,Q) - S_*(f,Q) < \frac \eps 2.
\]
The set $R=P\cup Q$ turns out to be a partition of $[a,b]$ on which we have
\begin{equation}\label{eq:56632}
S^*(f,R) = S^*(f,P) + S^*(f,Q), \qquad
S_*(f,R) = S_*(f,P) + S_*(f,Q)
\end{equation}
and thus
\[
S^*(f,R) - S_*(f,R) \le \frac \eps 2 + \frac \eps 2 = \eps.
\]
Applying the integrability criterion in reverse, we thus obtain
the integrability of $f$ on $[a,b]$ and the equations
\eqref{eq:56632} ensure the additivity of the integral with respect to the domain.

Conversely, if $f$ is integrable on $[a,b]$, the integrability criterion
guarantees that for every $\eps>0$ there exists a partition $R$ of $[a,b]$ such
that
\[
S^*(f,R) - S_*(f,R) < \eps.
\]
If we now consider $R' = R \cup \ENCLOSE{c}$, we know that $S^*(f,R') \le
S^*(f,R)$
and $S_*(f,R') \ge S_*(f,R)$, thus $R'$ also satisfies the property
\[
S^*(f,R') - S_*(f,R') < \eps.
\]
But now it is clear that setting $P=R'\cap[a,c]$ and $Q=R'\cap[c,b]$, it follows
that $P$
and $Q$ are partitions of $[a,c]$ and $[c,b]$ respectively, and that
\begin{align*}
  S^*(f,R') &= S^*(f,P) + S^*(f,Q), \\
  S_*(f,R') &= S_*(f,P) + S_*(f,Q).
\end{align*}
Thus we have
\begin{align*}
(S^*(f,P) - S_*(f,P)) + (S^*(f,Q) - S_*(f,Q))
&= S^*(f,R') - S_*(f,R')
< \eps.
\end{align*}
Since both addends $S^*-S_*$ are non-negative,
it follows that the inequalities
\[
S^*(f,P) - S_*(f,P) < \eps, \qquad
S^*(f,Q) - S_*(f,Q) < \eps.
\]
hold separately.
Thus $f$ is integrable both on $[a,c]$ and on $[c,b]$.
And once again, we can observe that the integral is additive over the domain.
\end{proof}  

\begin{theorem}[monotonicity of the integral]%
\label{th:monotonia_integrale}%
\mymark{*}%
Let $a\le b$ and let
$f,g\colon [a,b]\to \RR$ be two Riemann-integrable functions.
If for every $x\in [a,b]$ we have $f(x) \le g(x)$, then
\[
  \int_a^b f(x) \le \int_a^b g(x).
\]

In particular, if $f\ge 0$, then $\int_a^b f \ge 0$.
\end{theorem}
%
\begin{proof}
Clearly, if $f \le g$, then the $\sup$ of $f$ over any interval will be less
than or equal to the $\sup$ of $g$ over the same interval. Thus, on any
partition $P$ of $[a,b]$, we have:
\[
  S^*(f,P) \le S^*(g,P)
\]
from which we immediately obtain $I^*(f) \le I^*(g)$ and the result follows.
\end{proof}

\begin{theorem}[linearity of the integral]
\mymark{*}
\label{th:integrale_lineare}
Let $f,g\colon [a,b]\to \RR$ be two Riemann-integrable functions
and let $\lambda, \mu \in \RR$. Then $\lambda f + \mu g$
is Riemann-integrable and we have
\[
  \int_a^b (\lambda f + \mu g) = \lambda \int_a^b f + \mu \int_a^b g.
\]

In particular, the set of Riemann-integrable functions on $[a,b]$ is
a real vector space, and the integral is a
linear map on this space, taking values in $\RR$.
\end{theorem}
%
\begin{proof}
\mymark{*}
First, we show that
\begin{equation}\label{eq:20043}
  \int_a^b (-f) = -\int_a^b f.
\end{equation}
This follows from the fact that on any set $A$, we have
$\sup_A (-f) = -\inf_A f$, and thus for any partition $P$
we have
\[
  S^*(-f,P) = -S_*(f,P).
\]
It follows that $I^*(-f) = -I_*(f)$ and, similarly, $I_*(-f) = -I^*(f)$.
Thus, if $f$ is Riemann-integrable, then $-f$ is also integrable, and property
\eqref{eq:20043} holds.

Now, if $\lambda \ge 0$, we want to show that
\begin{equation}\label{eq:10032}
  \int_a^b \lambda f = \lambda \int_a^b f.
\end{equation}
Simply observe that $\sup_I \lambda f = \lambda \sup_I f$, and thus
$S^*(\lambda f,P) = \lambda S^*(f,P)$ for every partition $P$.
It follows that $I^*(\lambda f) = \lambda I^*(f)$. Similarly, we can
show that $I_*(\lambda f) = \lambda I_*(f)$. Thus, if $f$ is
Riemann-integrable, then $\lambda f$ (with $\lambda \ge 0$) is also integrable,
and \eqref{eq:10032} holds.

Combining \eqref{eq:20043} and $\eqref{eq:10032}$, we obtain
that $\eqref{eq:10032}$ holds for every $\lambda \in \RR$.
The same will be true if we replace $f$ with $g$ and $\lambda$ with $\mu$.
To conclude the proof, it will be sufficient
to show that
\begin{equation}\label{eq:80003}
\int_a^b (f+g) = \int_a^b f + \int_a^b g.
\end{equation}
Observe that on any set $A$, we have
\[
  \sup_A (f+g) \le \sup_A f + \sup_A g.
\]
Indeed, for every $x\in A$, we have $\sup_A f\ge f(x)$ 
and $\sup_A g \ge g(x)$, thus $\sup_A f + \sup_A g$
is greater than $f(x)+g(x)$ for every $x\in A$ 
and is therefore not less than $\sup_A (f+g)$, which 
is the least upper bound.
This means that
\[
  S^*(f+g) \le S^*(f) + S^*(g).
\]
Similarly, we can show that
\[
  S_*(f+g) \ge S_*(f) + S_*(g).
\]
Thus, we obtain
\[
  I^*(f+g) \le I^*(f) + I^*(g)
  \qquad\text{and}\qquad
  I_*(f+g) \ge I_*(f) + I_*(g)
\]
and thus, if $f$ and $g$ are integrable, then $f+g$ is also integrable,
and \eqref{eq:80003} holds.

To conclude that the set of integrable functions is a vector space,
it is sufficient to observe that, thanks to theorem~\ref{th:integrale_costante},
the function $0$ is integrable.
\end{proof}

\begin{definition}\label{def:osc}
Let $A\subset \RR$. We define 
the \emph{oscillation} of $A$ as:
\[
  \osc A \defeq \sup A - \inf A.
\]
In the case where $A=f(D)$ is the image of a function, 
we define the \emph{oscillation} of $f$ over the set 
$D$ as the oscillation of its image, i.e.:
\[
   \osc_{x\in D} f(x) 
   \defeq \osc f(D) 
   = \sup_{x\in D} f(x) - \inf_{x\in D} f(x).  
\]
\end{definition}

\begin{theorem}[integrability of continuous functions]
  \mymark{***}%
  \label{th:integrabilita_continue}%
  \mymargin{integrability of continuous functions}%
  \index{integrability!continuous functions}%
  Let $f\colon [a,b]\to \RR$ be a continuous function.
  Then $f$ is bounded and Riemann-integrable.
  \end{theorem}
  %
  \begin{proof}
  \mymark{***}
  By the Weierstrass theorem, we know that $f$ is bounded.
  By the Heine-Cantor theorem, we know that $f$ is uniformly continuous,
  thus for every $\eps>0$, there exists a $\delta>0$ such that
  \[
   \abs{x-y} < \delta \implies \abs{f(x)-f(y)} < \eps.
  \]
  We can then consider a partition $P_\delta$ with the property that
  the intervals determined by the partition all have a width less than
  $\delta$ (for example, we could take the partition formed by
    $(b-a)/\delta+2$ equally spaced points in $[a,b]$). On each interval $I$ of
  such
    a partition, we will have that if $x,y\in I$, then $\abs{f(x)-f(y)} < \eps$,
  from which
  it follows that $\sup_I f - \inf_I f \le \eps$.
  In particular, summing over all the intervals, we will have
  \begin{align*}
    S^*(f,P_\delta) - S_*(f,P_\delta)
        &= \sum_{k=1}^N (x_k-x_{k-1})\enclose{\sup_{[x_{k-1},x_k]} f -
    \inf_{[x_{k-1},x_k]} f} \\
    &\le \eps \sum_{k=1}^N (x_k - x_{k-1})
     = \eps (b-a).
  \end{align*}
  Since this quantity can be made arbitrarily small as
    $\eps \to 0$, based on the criteria for integrability, we can conclude that
  the
  function $f$ is integrable.
  \end{proof}
  
  \begin{theorem}[integrability of monotone functions]
  \label{th:integrabilita_monotone}%
  \mymargin{integrability of monotone functions}%
  \index{integrability!monotone functions}%
  Let $f\colon [a,b]\to \RR$ be a monotone function. Then $f$ is bounded and
  Riemann-integrable.
  \end{theorem}
  %
  \begin{proof}
  Suppose, for the sake of argument, that $f$ is increasing.
  
  Clearly, $f$ is bounded since $f(a) \le f(x) \le f(b)$ for every
  $x\in [a,b]$.
  
  To establish integrability, it suffices to show
  that there exists a sequence of partitions $P_n$
  such that $S^*(f,P_n) - S_*(f,P_n) \to 0$.
  Consider the equally spaced partition
  $P_n=\ENCLOSE{x_k \colon k=0,1, \dots, n}$ with $x_k=a+k(b-a)/N$.
  In this case, on each subinterval $[x_{k-1},x_k]$, we have
  \[
    \sup f([x_{k-1}, x_k]) = f(x_k),
    \qquad
    \inf f([x_{k-1}, x_k]) = f(x_{k-1}).
  \]
  Thus, the difference between the upper sums
  and the lower sums is telescopic
  and we have, as $n\to +\infty$
  \begin{align*}
  S^*(f,P) - S_*(f,P)
  &= \sum_{k=1}^n \frac{b-a}{n} f(x_k)
    - \sum_{k=1}^n \frac{b-a}{n} f(x_{k-1}) \\
  &= \frac{b-a}{n}(f(b)-f(a)) \to 0.
  \end{align*}
  This is what we wanted to demonstrate.
  \end{proof}
  
  \begin{example}[Heaviside function]
  \label{ex:heaviside}
  \mymargin{Heaviside function}
  \index{function!Heaviside}
  Let $a<0<b$.
  The function $H\colon [a,b] \to \RR$ defined by
  \[
  H(x) =
  \begin{cases}
  1 & \text{if $x\ge 0$}\\
  0 & \text{if $x< 0$}
  \end{cases}
  \]
  is increasing and therefore integrable.
  \end{example}
  
\begin{theorem}[integrability of the composite function]
  \label{th:composta_integrabile}%
  \index{integrability!of the composite function}%
  \index{composition!integrability}%
    Let $\phi\colon [\alpha,\beta]\to \RR$ be a continuous function and let
  $f\colon [a,b]\to [\alpha,\beta]$ be a Riemann-integrable function.
  Then the composite function 
  $\phi\circ f\colon[a,b]\to \RR$ is 
  bounded and Riemann-integrable.
\end{theorem}
%
\begin{proof}
According to the criteria for integrability,
theorem~\ref{th:criteri_integrabilita},
given $\eps>0$, it suffices to find a partition $P$
such that $S^*(\phi\circ f,P) - S_*(\phi\circ f,P)\le \eps(b-a)$.
By the Heine-Cantor theorem~\ref{th:heine_cantor}, 
the function $\phi$ is uniformly continuous on $[\alpha,\beta]$, 
and therefore there exists $\delta>0$ such that if $\abs{y_1-y_2}< \delta$,
then $\abs{\phi(y_1)-\phi(y_2)}< \eps$. 
By the Weierstrass theorem~\ref{th:weierstrass}, 
there exists $M$ such that $\abs{\phi(y)}\le M$ for every $y\in [\alpha,\beta]$.
In particular, this ensures that the composite function $\phi \circ f$ 
is bounded: $\phi(f(x))\le M$ for every $x\in [a,b]$.
Choosing $\sigma = \frac{\eps \delta}{4M} (b-a)$ 
(we will see at the end why this choice is made),
for the integrability of $f$, we know that 
there exists a partition $P$ consisting 
of the points $x_0=a < x_1 < \dots < x_n=b$ such that 
\mynote{The operator $\osc$ was introduced 
in definition~\ref{def:osc}}%
\[
  S^*(f,P) - S_*(f,P) 
  = \sum_{k=1}^n (x_k - x_{k-1})
  \cdot \osc f([x_{k-1},x_k])
  \le \sigma.
\]

We want to show that the following quantity is small:
\[
 S^*(\phi\circ f,P) - S_*(\phi\circ f,P)
 = \sum_{k=1}^n (x_k - x_{k-1}) \cdot \osc \phi(f([x_{k-1},x_k])).
\]
If $\osc f([x_{k-1},x_k]) \le \frac{\delta }2$, 
then for every $t,s\in[x_{k-1},x_k]$, we have 
$\abs{f(t)-f(s)}\le \frac{\delta}{2}$,
and therefore, by the uniform continuity of $\phi$,
$\abs{\phi(f(t))-\phi(f(s))}<\eps$.
Thus 
\[
  \osc f([x_{k-1},x_k]) \le \frac{\delta}{2} 
  \implies  
  \osc \phi(f([x_{k-1},x_k]))\le \eps.
\] 
On the other hand, the intervals for which 
$\osc f([x_{k-1},x_k]) > \frac{\delta}{2}$ 
must be few, otherwise the function would not be 
integrable.
If we call $K$ the set of indices $k$ for which 
$M_k-m_k\ge \frac \delta 2$, it must be 
\[
  \frac{\delta}{2}\sum_{k\in K} (x_k-x_{k-1})
  \le \sum_{k\in K} (x_k-x_{k-1})\cdot \osc f([x_{k-1},x_k])
  \le S^*(f,P) - S_*(f,P) \le \sigma.
\]
Since the function $\phi$ takes values between $-M$ and $M$, 
we have  
\[
   \sum_{k\in K} (x_k-x_{k-1})
   \osc \phi(f([x_{k-1},x_k]))
   \le 2M \sum_{k\in K} (x_k-x_{k-1})
   \le \frac{4M\sigma}{\delta} = \eps (b-a).
\]

In conclusion: 
\begin{align*}
  S^*(\phi\circ f,P)-S_*(\phi\circ f,P)
  \le 
  \eps (b-a) + 
  \sum_{k\not \in K} (x_k-x_{k-1}) \eps
  = 2\eps (b-a)
\end{align*}
which, according to the criteria for integrability, gives us the thesis.
\end{proof}

\begin{remark}
If in the previous theorem we take $f(x)=x$, we obtain 
the theorem~\ref{th:integrabilita_continue} on the integrability of continuous
functions.
\end{remark}

Given $a,b\in \RR$, we define:
\begin{align*}
  a^+ &= \begin{cases}
    a & \text{if $a>0$,}\\
    0 & \text{otherwise,}
  \end{cases} &
  a^- &= \begin{cases}
    -a & \text{if $a<0$,}\\
    0 & \text{otherwise,}
  \end{cases} \\
  a \vee b &= \begin{cases}
    a & \text{if $a\ge b$,}\\
    b & \text{otherwise,}
  \end{cases} &
  a \wedge b &= \begin{cases}
    a & \text{if $a\le b$,}\\
    b & \text{otherwise.}
  \end{cases}
\end{align*}
The values $a^+$ and $a^-$ are called the 
\emph{positive part} and \emph{negative part} of $a$, respectively.
The operation $a\vee b = \max\ENCLOSE{a,b}$ returns the maximum between $a$ and
$b$,
while $a\wedge b = \min\ENCLOSE{a,b}$ returns the minimum between $a$ and $b$.
Note that the following properties hold:
\[
   a^+ = a \vee 0, \qquad
   a^- = -(a \wedge 0), \qquad
   a = a^+ - a^-, \qquad
   \abs{a} = a^+ + a^-.
\]
Moreover, it is easy to convince oneself that the following holds:
\[
  a + b = (a \vee b) + (a \wedge b), \qquad
  \abs{a-b} = (a \vee b) - (a \wedge b).
\]

If $f$ and $g$ are real-valued functions, we define 
the functions $f^+$, $f^-$, $f\vee g$, and $f\wedge g$ accordingly:
\begin{gather*}
  f^+(x) = (f(x))^+, \qquad
  f^-(x) = (f(x))^-, \\
  (f\vee g)(x) = f(x) \vee g(x), \qquad
  (f\wedge g)(x) = f(x) \wedge g(x).
\end{gather*}

\begin{theorem}[lattice properties]%
\label{th:reticolo}%
If $f$ is a Riemann-integrable function on the interval $[a,b]$, then
$\abs{f}$, $f^+$, and $f^-$ are also integrable, and if
$g$ is also Riemann-integrable on $[a,b]$, then $f\vee g$ and $f\wedge g$
are integrable on the same interval.
\end{theorem}
%
\begin{proof}
  It suffices to observe that the functions 
  $y\mapsto \abs{y}$, $y\mapsto y^+$, and $y\mapsto y^-$ 
  are continuous functions, and apply theorem~\ref{th:composta_integrabile}
  on the integrability of the composite function.
  \end{proof}
  %
  \begin{proof}[Alternative proof.]
  We can provide a proof that does not require 
  theorem~\ref{th:composta_integrabile}
  but essentially follows its proof in the 
  much simpler case that interests us.

    First, we demonstrate that if $f$ is integrable, then $f^+$ is also
  integrable.
  Since $\abs{f^+(x)-f^+(y)}\le \abs{f(x)-f(y)}$, 
  it is easy to convince oneself that on any set $A$, it holds that 
  \mynote{The operator $\osc$ was introduced 
  in definition~\ref{def:osc}}%
  $\osc_A f^+ \le \osc_A f$.
  Thus, on any partition $P$,
  we obtain:
  \[
    0 \le S^*(f^+,P) - S_*(f^+,P) \le S^*(f,P) - S_*(f,P).
  \]
  By the first criterion for integrability in 
  theorem~\ref{th:criteri_integrabilita}, since $f$ is integrable, 
  we know that for every $\eps>0$, there exists $P$ such that the right-hand 
  side of the previous equation is less than $\eps$.
  But then the left-hand side is also less than $\eps$, and thus, 
  by the same criterion, $f^+$ is also integrable.    
  Similarly, it can be shown that if $f$ is integrable, 
  then $f^-$ is integrable.
  \mynote{
    Alternatively, one can observe that $f^- = -(-f)^+$ and exploit the 
    fact that if $f$ is integrable, we already know that $-f$ 
    is integrable.
  }
  
  Regarding $\abs f$, $f\vee g$, and $f\wedge g$,  
  it suffices to observe that%
  \mynote{For every $a,b\in \RR$, 
  it is easy to convince oneself that the following identities hold:
  \[
    \begin{cases}
    a = a^+ - a^-\\
    \abs{a} = a^+ + a^-\\
    \end{cases}
  \]\[
    \begin{cases}
      a + b = (a\vee b) + (a\wedge b)\\
      \abs{a - b} = (a\vee b) - (a\wedge b) 
    \end{cases}
  \]}
  \begin{gather*}
    \abs{f} = f^+ + f^-, \qquad
  %  f^+ = \frac{\abs{f} + f}{2}, \qquad
  %  f^- = \frac{\abs{f} - f}{2}, \\
    f\wedge g = \frac{f + g - \abs{f-g}}{2}, \qquad
    f\vee g = \frac{f+g + \abs{f-g}}{2}
  \end{gather*}
  and since linear combinations of integrable functions 
  are integrable functions (theorem~\ref{th:integrale_lineare}),
  the integrability of the functions on the left-hand side
  of the equalities is obtained.
\end{proof}

\begin{remark}
Observe that if $f$ and $g$ are continuous functions, 
then $f^+$, $f^-$, $\abs f$, $f\vee g$, and $f\wedge g$ are also continuous
functions
and thus Riemann integrable, thanks to theorem~\ref{th:integrabilita_continue}.
Therefore, the previous theorem is only useful when dealing with 
Riemann integrable functions that are not known to be continuous. 
\end{remark}

\begin{theorem}
\label{th:integrale_modulo}%
If $a\le b$ and $f$ is a bounded and Riemann-integrable function on $[a,b]$,
then 
\[
  \abs{\int_a^b f(x)\, dx} \le \int_a^b \abs{f(x)}\, dx.
\]
In general (without the assumption $a \le b$), we have
\begin{equation}\label{eq:489732}
  \abs{\int_a^b f(x)\, dx} 
  \le 
  \abs{\int_a^b \abs{f(x)}}. 
\end{equation}
\end{theorem}
\begin{proof}
Observe that the functions $\abs{f}$, $f^+$, and $f^-$ are Riemann-integrable
by the previous theorem~\ref{th:reticolo}.

If $a\le b$, since $f^+\ge 0$ and $f^-\ge 0$,
by monotonicity (theorem~\ref{th:monotonia_integrale}),
we know that $\int_a^b f^+\ge 0$ 
and $\int_a^b f^-\ge 0$, so that, 
by linearity (theorem~\ref{th:integrale_lineare}):
\begin{align*}
\abs{\int_a^b f} 
&= \abs{\int_a^b f^+ - \int_a^b f^-}\\
&\le \abs{\int_a^b f^+} + \abs{\int_a^b f^-}
= \int_a^b f^+ + \int_a^b f^-
= \int_a^b \abs{f}.   
\end{align*}
If $b\le a$, we can reduce to the previous case by swapping 
the limits of integration.
\end{proof}

The following result is a continuous version of the convexity inequality 
for barycentric combinations.
It can be noted that the previous theorem~\ref{th:integrale_modulo} is a special
case
of the following with $\phi(x)=\abs{x}$.
%
\begin{theorem}[Jensen's inequality]
  \label{th:jensen}%
  \index{inequality!Jensen's}%
  \index{Jensen!inequality}%
  \index{theorem!Jensen}%
Let $f\colon[a,b]\to I\subset \RR$ be a continuous function taking values 
in an interval $I$, and let $\phi\colon I \to \RR$ be a convex and continuous
function.
Let $g\colon [a,b]\to \RR$ be a continuous,
non-negative function such that 
\[
    \int_a^b g(x)\, dx = 1
\]
Then we have 
  \[
  \phi\enclose{\int_a^b f(x) g(x) \, dx} \le \int_a^b \phi(f(x)) g(x)\, dx.
  \]
In the particular case where $g$ is chosen as the constant function
$g(x)=1/(b-a)$, we obtain
the more expressive formulation:
\[
  \phi\enclose{\dashint_a^b f(x)\, dx} \le \dashint_a^b \phi(f(x))\, dx
\]
where $\dashint$ denotes the integral mean:
\[
  \dashint_a^b f(x)\, dx = \frac{1}{b-a} \int_a^b f(x)\, dx.
\]
\end{theorem}
%
\begin{proof}
The continuity assumptions on $f$, $g$, and $\phi$ are redundant, but they
provide
an immediate guarantee that both integrand functions
are indeed Riemann-integrable.
Let 
\[
  y_0 = \int_a^b f(x) g(x)\, dx 
\]
choose an affine function $L(y) = my + q$ such that $L(y_0)=\phi(y_0)$ 
and $L(y)\le \phi(y)$ for every $y\in I$ (thanks to
theorem~\ref{th:supporto_convessa}).
By the linearity of the integral, we have: 
\[
  L(y_0)
  = \int_a^b mf(x)g(x)\, dx + q \int_a^b g(x)\, dx
  = \int_a^b L(f(x)) g(x)\, dx.
\]
Then, by the monotonicity of the integral, it follows that
\[
  \phi(y_0) = L(y_0) = \int_a^b L(f(x)) g(x)\, dx \le \int_a^b \phi(f(x)) g(x)\, dx.
\]
\end{proof}
    

\begin{theorem}[continuity of the integral]
\label{th:integrale_continuo}
Let $f\colon [a,b]\to \RR$ be a bounded function such
that for every $c\in (a,b]$, $f$ is Riemann-integrable
on $[c,b]$. Then $f$ is Riemann-integrable on $[a,b]$ and it holds that
\[
  \int_a^b f = \lim_{c\to a^+} \int_c^b f.
\]

Similarly, if $f$ is integrable on every interval $[a,c]$
with $c\in [a,b)$, then $f$ is integrable on $[a,b]$ and it holds that
\[
  \int_a^b f = \lim_{c\to b^-} \int_a^c f.
\]
\end{theorem}
%
\begin{proof}
We will only prove the first part, as the second is treated
in an entirely analogous manner. Suppose, therefore, that $f$ is integrable on
every
interval $[c,b]$ with $c\in (a,b]$. We also know that $f$ is bounded on
the entire interval $[a,b]$, so there exists $M>0$ such that $\abs{f(x)}\le M$
for every
$x\in [a,b]$. Given any $\eps>0$,
choose $c = a + \eps/(4M)$ and, knowing that $f$ is integrable on $[c,b]$,
consider a partition $Q_\eps$ of $[c,b]$ such that
\[
  S^*(f, Q_\eps) - S_*(f, Q_\eps) < \frac \eps 2.
\]
Then, setting $P_\eps = \ENCLOSE{a} \cup Q_\eps$, we obtain:
\begin{align*}
  S^*(f,P_\eps) &= S^*(f,Q_\eps) + (c-a)\sup_{[a,c]} f
    \le S^*(f,Q_\eps) + M \cdot (c-a) \le S^*(f,Q_\eps) + \frac \eps 2\\
  S_*(f,P_\eps) &= S_*(f,Q_\eps) + (c-a)\inf_{[a,c]} f
    \ge S_*(f,Q_\eps) - M \cdot (c-a) \ge S_*(f,Q_\eps) - \frac \eps 2
\end{align*}
from which
\begin{align*}
  S^*(f,P_\eps)-S_*(f,P_\eps)
  &\le S^*(f,Q_\eps) - S_*(f,Q\eps) + 2 M \cdot (c-a)\\
  &< \frac{\eps}{2} + 2 M \frac{\eps}{4M}
  = \eps.
\end{align*}
Thus, the function $f$ is integrable. But we have
\begin{align*}
\int_a^b f \le S^*(f,P_\eps) \le S^*(f,Q_\eps) + \frac{\eps}{4}
\le S_*(f,Q_\eps) + \frac 3 4 \eps
\le \int_{a+\frac {\eps}{4M}}^b f + \frac {3}{2} \eps \\
\int_a^b f \ge S_*(f,P_\eps) \ge S_*(f,Q_\eps) - \frac{\eps}{4}
\ge S^*(f,Q_\eps) - \frac 3 4 \eps
\ge \int_{a+\frac {\eps}{4M}}^b f - \frac {3}{2} \eps
\end{align*}
from which, taking the limit as $\eps \to 0^+$, we obtain
\[
  \lim_{c\to a^+} \int_c^b f \le \int_a^b f \le \lim_{c\to a^+} \int_c^b f
\]
and therefore the equality.
\end{proof}

\begin{example}
The function
\[
  f(x) = \begin{cases}
  \sin\frac 1 x & \text{if $x\neq 0$,}\\
  0 & \text{if $x=0$}
  \end{cases}
\]
is integrable on every closed and bounded interval $[a,b]$.
\end{example}
\begin{proof}
On every interval $[\eps,b]$ with $\eps>0$ and $b>\eps$, the function
is integrable since it is continuous on such intervals (the function is
continuous on
the entire $\RR\setminus\ENCLOSE{0}$). Moreover, the function is bounded, so by
the previous theorem
we can conclude that it is integrable on $[0,b]$ for every $b>0$.

Similarly (by symmetry), the function is integrable on $[a,0]$
for every $a<0$.
Thus, by additivity, the function is integrable on every $[a,b]$ with $a<0$ and
$b>0$,
and consequently (again thanks to theorem~\ref{th:additivita_integrale})
it is integrable on every interval $[a,b]$.
\end{proof}